\section{Limitations}

\benchmark{} is designed as a controlled first step toward evaluating
accounting reliability under exact ground truth.
Its problems are procedurally generated rather than drawn from live
filings, which enables deterministic scoring while still modeling the
core workflow of balance-sheet preparation from opening balances,
journal entries, adjustments, and closing behavior.
Future benchmark extensions can broaden the distribution of language and
disclosure styles through de-identified real-world statements,
expert-authored variants, or hybrid templates layered on the same ledger
engine.

The empirical study also uses two evaluation pools---Large-500 and
Controlled-100---to balance statistical precision with frontier-model
coverage.
All inferential claims are restricted to matched comparisons within a
pool, and future work can consolidate more systems onto a single
fully matched leaderboard as additional large-pool runs become
available.

Our current focus is statement construction rather than downstream
auditing or interactive financial analysis.
That decomposition is intentional: statement construction isolates the
core bookkeeping problem before introducing broader audit judgments.
Future work can extend the benchmark toward audit-style tasks such as
error localization, correction, and explanation.

Finally, all experiments use end-to-end generation without external
calculation or ledger tools.
This setting exposes what the models can do unaided, but tool-augmented
variants would be useful for separating accounting classification errors
from arithmetic execution failures.
